\documentclass[12pt,a4paper]{report}
\usepackage[bahasa]{babel}
\usepackage[utf8]{inputenc}
\usepackage[T1]{fontenc}
\usepackage{geometry}
\geometry{a4paper, margin=2.5cm, top=3cm, bottom=2.5cm, headheight=14.5pt}
\usepackage{xcolor}
\usepackage[most]{tcolorbox}
\usepackage{graphicx}
\usepackage{booktabs}
\usepackage{tabularx}
\usepackage{enumitem}
\usepackage[expansion=false]{microtype}
\usepackage{titlesec}
\usepackage{fancyhdr}
\usepackage{amsmath}
\usepackage{amssymb}

% Warna Korporat PLN
\definecolor{plnblue}{RGB}{0,75,135}
\definecolor{plnorange}{RGB}{255,107,53}
\definecolor{fillgreen}{RGB}{240,255,240}

\usepackage{hyperref}
\hypersetup{
    colorlinks=true,
    linkcolor=plnblue,
    urlcolor=plnblue,
    citecolor=plnblue
}

% Header & Footer
\pagestyle{fancy}
\fancyhf{}
\fancyhead[L]{\small\textbf{Post-Test Modul 4.3}}
\fancyhead[R]{\small Level 4 (Mastery) -- PLN}
\fancyfoot[C]{\small\thepage}
\renewcommand{\headrulewidth}{0.4pt}

% Format Heading
\titleformat{\chapter}[display]{\normalfont\huge\bfseries\color{plnblue}}{\chaptertitlename\ \thechapter}{20pt}{\Huge}
\titleformat{\section}{\Large\bfseries\color{plnblue}}{}{0em}{}
\titleformat{\subsection}{\large\bfseries\color{plnblue}}{}{0em}{}
\titlespacing*{\chapter}{0pt}{-30pt}{20pt}
\titlespacing*{\section}{0pt}{12pt}{8pt}
\titlespacing*{\subsection}{0pt}{10pt}{6pt}

% Custom Boxes
\newtcolorbox{plnbox}[1][]{%
  colback=plnblue!5!white, colframe=plnblue, fonttitle=\bfseries,
  title=#1, sharp corners, boxrule=0.8pt, breakable
}
\newtcolorbox{warnbox}{%
  colback=plnorange!10, colframe=plnorange, fonttitle=\bfseries,
  title=Petunjuk Penting, sharp corners, breakable
}
\newtcolorbox{answerbox}{%
  colback=fillgreen, colframe=green!50!black, fonttitle=\bfseries,
  title=Kunci Jawaban (Khusus Fasilitator), sharp corners, breakable
}

% List styling
\setlist[itemize]{leftmargin=*, topsep=2pt, itemsep=2pt}
\setlist[enumerate]{leftmargin=*, topsep=2pt, itemsep=2pt}

% Helper: answer lines
\newcommand{\ansline}[1][5.5cm]{\underline{\hspace{#1}}}

\begin{document}

% ================= COVER =================
\begin{titlepage}
\centering
\vspace*{1.5cm}
{\Huge\bfseries\color{plnblue} POST-TEST}\\[0.4cm]
{\LARGE\bfseries Modul 4.3: Optimalisasi Prosedur dan}\\[0.2cm]
{\LARGE\bfseries Alur Kerja Terintegrasi}\\[0.8cm]
\rule{0.8\linewidth}{0.6pt}\\[0.6cm]
{\large Program Penguatan Kompetensi Tata Kelola Kearsipan}\\
{\large Level 4 (Mastery)}\\[1.2cm]
{\large Disusun Khusus untuk: \textbf{PT PLN (Persero)}}\\
{\large Target Peserta: Manajemen Atas (MA)}\\[2.5cm]
\begin{figure}[h]
\centering
\includegraphics[width=4cm]{example-image}
\end{figure}
\vfill
{\small Waktu Pengerjaan: 30 Menit \quad|\quad Nilai Kelulusan Minimal: 65/100}
\end{titlepage}

% ================= IDENTITAS PESERTA =================
\newpage
\thispagestyle{fancy}

\begin{center}
{\LARGE\bfseries\color{plnblue} LEMBAR JAWABAN POST-TEST}\\[0.3cm]
{\large Modul 4.3 -- Level 4 (Mastery)}
\end{center}

\vspace{0.6cm}

\renewcommand{\arraystretch}{1.8}
\noindent
\begin{tabularx}{\textwidth}{@{} X X @{}}
\textbf{Nama Peserta:} \ansline & \textbf{Unit Kerja:} \ansline \\
\textbf{Jabatan:} \ansline & \textbf{Tanggal:} \ansline \\
\end{tabularx}
\renewcommand{\arraystretch}{1.0}

\vspace{0.4cm}

\begin{warnbox}
\begin{itemize}[nosep]
    \item Waktu pengerjaan: \textbf{30 menit}. Kerjakan seluruh bagian (A--D).
    \item Dilarang membuka modul atau catatan selama pengerjaan.
    \item Untuk Bagian A \& B, lingkari atau beri tanda \textbf{X} pada jawaban yang dipilih.
    \item Untuk Bagian C \& D, tulis jawaban pada ruang yang disediakan. Gunakan bahasa yang ringkas dan substansial.
    \item Nilai kelulusan minimal: \textbf{65 dari 100}.
\end{itemize}
\end{warnbox}

\vspace{0.3cm}

\renewcommand{\arraystretch}{1.3}
\begin{center}
\small
\begin{tabularx}{0.85\textwidth}{@{} X c c c @{}}
\toprule
\textbf{Bagian} & \textbf{Jumlah Soal} & \textbf{Bobot/Soal} & \textbf{Total Poin} \\
\midrule
A. Pilihan Ganda & 10 soal & 3 poin & 30 poin \\
B. Benar/Salah & 5 soal & 2 poin & 10 poin \\
C. Jawaban Singkat & 5 soal & 6 poin & 30 poin \\
D. Studi Kasus Analitis & 2 soal & 15 poin & 30 poin \\
\midrule
\textbf{TOTAL} & \textbf{22 soal} & & \textbf{100 poin} \\
\bottomrule
\end{tabularx}
\end{center}
\renewcommand{\arraystretch}{1.0}

% ================= BAGIAN A: PILIHAN GANDA =================
\newpage
\section{BAGIAN A: Pilihan Ganda (10 soal $\times$ 3 poin = 30 poin)}
\textit{Lingkari SATU jawaban yang paling tepat untuk setiap pertanyaan.}

\vspace{0.3cm}

\begin{enumerate}[leftmargin=*, itemsep=0.5cm]

% Soal 1 -- Kriteria 4.3.1 (Bab 1-2: Definisi Optimalisasi)
\item Dalam konteks modul ini, ``optimalisasi prosedur'' kearsipan pada lingkungan \textit{Arsip Hibrida} berarti:
\begin{enumerate}[label=\Alph*., nosep]
    \item Menghapus seluruh dokumen fisik dan memaksa 100\% \textit{paperless}
    \item Mengeliminasi langkah non-nilai tambah dan mengotomasi \textit{metadata tracking}, meskipun sebagian dokumen tetap fisik
    \item Membeli perangkat lunak manajemen dokumen terbaru dari vendor ternama
    \item Menambah jumlah staf kearsipan untuk mempercepat pemrosesan manual
\end{enumerate}

% Soal 2 -- Kriteria 4.3.1 (Bab 2: Prinsip Evaluasi PARBICA)
\item Prinsip \textit{Single Source of Truth} dalam evaluasi prosedur kearsipan mengharuskan bahwa:
\begin{enumerate}[label=\Alph*., nosep]
    \item Seluruh data arsip hanya boleh disimpan di satu server tunggal
    \item Data arsip tercipta sekali di titik transaksi bisnis dan mengalir konsisten tanpa input ulang manual
    \item Hanya satu orang yang boleh mengakses arsip pada satu waktu
    \item Semua arsip harus memiliki satu format file yang seragam (PDF)
\end{enumerate}

% Soal 3 -- Kriteria 4.3.1 (Bab 2: Metode Fact-Finding)
\item Teknik \textit{Document Shadowing} dalam pengumpulan fakta lapangan adalah:
\begin{enumerate}[label=\Alph*., nosep]
    \item Membuat salinan bayangan (\textit{shadow copy}) dari seluruh arsip digital
    \item Melacak satu dokumen nyata dari penciptaan hingga penyimpanan final, mencatat setiap hambatan dan waktu transisi
    \item Mengamati cara staf mengoperasikan komputer dari belakang bahu mereka
    \item Membandingkan isi dokumen asli dengan salinannya untuk mendeteksi pemalsuan
\end{enumerate}

% Soal 4 -- Kriteria 4.3.2 (Bab 3: Prinsip Desain Workflow)
\item Prinsip \textit{Rule-Based Routing} dalam desain alur kerja terintegrasi berarti:
\begin{enumerate}[label=\Alph*., nosep]
    \item Persetujuan dokumen diarahkan secara otomatis berdasarkan hierarki, nilai transaksi, atau klasifikasi risiko
    \item Seluruh dokumen harus melewati semua tingkat manajer tanpa pengecualian
    \item Routing dilakukan secara manual berdasarkan arahan lisan atasan
    \item Dokumen dikirim ke seluruh karyawan secara \textit{broadcast} untuk transparansi
\end{enumerate}

% Soal 5 -- Kriteria 4.3.2 (Bab 3: Logical Fallacy Parallel Routing)
\item Menurut perspektif arsiparis, risiko dari \textit{Parallel Digital Routing} tanpa mekanisme \textit{Immutability Lock} adalah:
\begin{enumerate}[label=\Alph*., nosep]
    \item Dokumen terlalu cepat disetujui sehingga tidak sempat diperiksa
    \item Terjadi ``bencana \textit{Version Control}'' -- arsiparis tidak tahu revisi mana yang menjadi \textit{Golden Record}
    \item Biaya bandwidth jaringan meningkat drastis
    \item Supervisor kehilangan wewenang karena keputusan diambil secara paralel
\end{enumerate}

% Soal 6 -- Kriteria 4.3.3 (Bab 4: SOP Digital)
\item Perbedaan mendasar antara ``SOP IT'' dan ``SOP Kearsipan'' menurut arsiparis adalah:
\begin{enumerate}[label=\Alph*., nosep]
    \item SOP IT menggunakan bahasa Inggris, SOP Kearsipan menggunakan bahasa Indonesia
    \item SOP Kearsipan wajib mencakup Klasifikasi, Retensi, dan Hak Akses (RBAC), bukan sekadar metadata teknis
    \item SOP IT dirancang oleh Divisi IT, SOP Kearsipan dirancang oleh Divisi Umum
    \item Keduanya sebenarnya identik, hanya berbeda nama
\end{enumerate}

% Soal 7 -- Kriteria 4.3.3 (Bab 4: Komponen SOP)
\item Berapa jumlah komponen wajib minimum yang harus dimuat dalam SOP Digital Terintegrasi menurut modul ini?
\begin{enumerate}[label=\Alph*., nosep]
    \item 5 komponen
    \item 6 komponen
    \item 7 komponen
    \item 8 komponen
\end{enumerate}

% Soal 8 -- Kriteria 4.3.3 (Bab 4: Exception Handling)
\item Jika Sistem Tata Persuratan Digital mengalami gangguan (\textit{down}) lebih dari 4 jam, langkah \textit{Exception Handling} yang benar menurut SOP adalah:
\begin{enumerate}[label=\Alph*., nosep]
    \item Menunggu sistem pulih tanpa melakukan apa pun
    \item Mengirim dokumen via WhatsApp sebagai pengganti sementara
    \item Menggunakan prosedur input manual terstandar, lalu menyinkronkan data setelah sistem pulih
    \item Mencetak seluruh dokumen dan memproses secara konvensional permanen
\end{enumerate}

% Soal 9 -- Kriteria 4.3.4 (Bab 5: KPI)
\item KPI ``\textit{Retrieval Accuracy Rate}'' dalam pengukuran kinerja kearsipan mengukur:
\begin{enumerate}[label=\Alph*., nosep]
    \item Kecepatan koneksi internet saat mengunduh arsip
    \item Persentase arsip yang berhasil ditemukan pada pencarian pertama tanpa revisi kueri
    \item Jumlah total arsip yang tersimpan di repositori
    \item Tingkat kepuasan pengguna terhadap tampilan antarmuka sistem
\end{enumerate}

% Soal 10 -- Kriteria 4.3.4 (Bab 5: Integrasi Kearsipan)
\item Menurut perspektif arsiparis, ``Integrasi Kearsipan'' yang sesungguhnya menuntut arsip bergerak utuh bersama tiga elemen. Elemen-elemen tersebut adalah:
\begin{enumerate}[label=\Alph*., nosep]
    \item Hardware, Software, dan Brainware
    \item Input, Proses, dan Output
    \item \textit{Content} (isi), \textit{Structure} (format), dan \textit{Context} (metadata pencipta \& kronologi)
    \item Kecepatan, Kapasitas, dan Keamanan
\end{enumerate}

\end{enumerate}

% ================= BAGIAN B: BENAR/SALAH =================
\newpage
\section{BAGIAN B: Benar/Salah (5 soal $\times$ 2 poin = 10 poin)}
\textit{Tulis \textbf{B} (Benar) atau \textbf{S} (Salah) pada kotak yang disediakan. Jika Salah, tuliskan koreksi singkat pada baris di bawahnya.}

\vspace{0.3cm}

\begin{enumerate}[leftmargin=*, itemsep=0.6cm]

\item $\square$ \quad Jika dua aplikasi berhasil mentransfer file PDF dari Server A ke Server B melalui API, maka kedua sistem tersebut sudah ``terintegrasi'' secara kearsipan.\\[0.2cm]
Koreksi (jika Salah): \ansline[12cm]

\item $\square$ \quad Pada pendekatan Arsip Hibrida, dokumen kontrak yang diwajibkan fisik oleh regulasi audit tetap perlu diotomasi \textit{metadata tracking}-nya di sistem digital.\\[0.2cm]
Koreksi (jika Salah): \ansline[12cm]

\item $\square$ \quad Evaluasi prosedur kearsipan dalam modul ini bertujuan untuk mengaudit konfigurasi teknis server dan infrastruktur jaringan.\\[0.2cm]
Koreksi (jika Salah): \ansline[12cm]

\item $\square$ \quad Komponen ``Klasifikasi, Retensi, \& Akses'' adalah komponen opsional dalam SOP Digital Terintegrasi dan boleh diabaikan jika sistem IT sudah mencatat metadata.\\[0.2cm]
Koreksi (jika Salah): \ansline[12cm]

\item $\square$ \quad Teknik \textit{Gemba Walk} dalam pengumpulan fakta lapangan dilakukan dengan mengamati dan mewawancarai langsung staf operasional pembuat arsip mengenai kendala harian mereka.\\[0.2cm]
Koreksi (jika Salah): \ansline[12cm]

\end{enumerate}

% ================= BAGIAN C: JAWABAN SINGKAT =================
\newpage
\section{BAGIAN C: Jawaban Singkat (5 soal $\times$ 6 poin = 30 poin)}
\textit{Jawab dengan ringkas dan substansial (maksimal 3--4 kalimat per soal).}

\vspace{0.3cm}

\begin{enumerate}[leftmargin=*, itemsep=0.3cm]

% Soal C1 -- Kriteria 4.3.1
\item Sebutkan dan jelaskan secara singkat \textbf{3 teknik pengumpulan fakta lapangan (\textit{Fact-Finding})} yang diajarkan di modul ini untuk menggali data prosedur yang ada.

\vspace{0.2cm}
\noindent\rule{\textwidth}{0.3pt}
\vspace{0.5cm}
\noindent\rule{\textwidth}{0.3pt}
\vspace{0.5cm}
\noindent\rule{\textwidth}{0.3pt}
\vspace{0.5cm}
\noindent\rule{\textwidth}{0.3pt}

% Soal C2 -- Kriteria 4.3.1
\item Dari 7 domain pada ``Checklist Diagnostik Kesiapan Digital'' yang diadaptasi dari PARBICA Guideline 13, sebutkan \textbf{2 domain} beserta masing-masing satu pertanyaan diagnostik yang paling relevan untuk unit kerja Anda.

\vspace{0.2cm}
\noindent\rule{\textwidth}{0.3pt}
\vspace{0.5cm}
\noindent\rule{\textwidth}{0.3pt}
\vspace{0.5cm}
\noindent\rule{\textwidth}{0.3pt}
\vspace{0.5cm}
\noindent\rule{\textwidth}{0.3pt}

% Soal C3 -- Kriteria 4.3.2
\item Jelaskan perbedaan antara pola integrasi \textbf{\textit{Trigger-Based Auto-Capture}} dan \textbf{\textit{SLA-Driven Escalation}} dalam konteks alur kerja kearsipan PLN.

\vspace{0.2cm}
\noindent\rule{\textwidth}{0.3pt}
\vspace{0.5cm}
\noindent\rule{\textwidth}{0.3pt}
\vspace{0.5cm}
\noindent\rule{\textwidth}{0.3pt}
\vspace{0.5cm}
\noindent\rule{\textwidth}{0.3pt}

% Soal C4 -- Kriteria 4.3.3
\item Mengapa komponen ``Klasifikasi, Retensi, \& Akses'' dijadikan komponen wajib ke-5 dalam struktur SOP Digital Terintegrasi? Jelaskan dampaknya jika komponen ini tidak ada.

\vspace{0.2cm}
\noindent\rule{\textwidth}{0.3pt}
\vspace{0.5cm}
\noindent\rule{\textwidth}{0.3pt}
\vspace{0.5cm}
\noindent\rule{\textwidth}{0.3pt}
\vspace{0.5cm}
\noindent\rule{\textwidth}{0.3pt}

% Soal C5 -- Kriteria 4.3.4
\item Sebutkan \textbf{3 komponen penghematan} dalam pendekatan \textit{Time-Cost-Risk Saving} untuk menghitung ROI optimalisasi prosedur kearsipan.

\vspace{0.2cm}
\noindent\rule{\textwidth}{0.3pt}
\vspace{0.5cm}
\noindent\rule{\textwidth}{0.3pt}
\vspace{0.5cm}
\noindent\rule{\textwidth}{0.3pt}
\vspace{0.5cm}
\noindent\rule{\textwidth}{0.3pt}

\end{enumerate}

% ================= BAGIAN D: STUDI KASUS =================
\newpage
\section{BAGIAN D: Studi Kasus Analitis (2 soal $\times$ 15 poin = 30 poin)}
\textit{Baca skenario dengan cermat, lalu jawab pertanyaan yang mengikutinya.}

\vspace{0.3cm}

% === STUDI KASUS 1 ===
\begin{plnbox}[Skenario 1: Unit Transmisi PLN Regional X]
\small
Unit Transmisi PLN Regional X mengelola arsip laporan inspeksi rutin untuk 32 menara transmisi 150 kV. Saat ini prosedurnya sebagai berikut:
\begin{enumerate}[nosep]
    \item Teknisi lapangan mengisi formulir inspeksi \textbf{kertas} di lokasi menara.
    \item Formulir diantar secara fisik ke kantor area (memakan waktu 1--2 hari).
    \item Admin area mengetik ulang data inspeksi ke dalam \textit{spreadsheet} Excel.
    \item Admin mencetak \textit{spreadsheet}, meminta tanda tangan basah Supervisor.
    \item Dokumen yang sudah ditandatangani di-scan, lalu diupload ke folder bersama di server lokal.
    \item Metadata (nama menara, tanggal inspeksi, kategori temuan) diisi manual oleh admin -- sering terjadi inkonsistensi penamaan file.
\end{enumerate}
Rata-rata \textit{cycle time} dari inspeksi hingga arsip tersimpan: \textbf{5--7 hari kerja}. Dalam 12 bulan terakhir, 23\% arsip inspeksi tidak dapat ditemukan saat dibutuhkan untuk audit BPK.
\end{plnbox}

\vspace{0.3cm}

\begin{enumerate}[leftmargin=*, label=\textbf{D1.\alph*}]
    \item \textbf{(5 poin)} Identifikasi minimal \textbf{3 titik gesekan (\textit{friction points})} dari prosedur di atas menggunakan prinsip \textit{Value-Added Focus}.

\vspace{0.2cm}
\noindent\rule{\textwidth}{0.3pt}
\vspace{0.5cm}
\noindent\rule{\textwidth}{0.3pt}
\vspace{0.5cm}
\noindent\rule{\textwidth}{0.3pt}
\vspace{0.5cm}
\noindent\rule{\textwidth}{0.3pt}
\vspace{0.5cm}
\noindent\rule{\textwidth}{0.3pt}

    \item \textbf{(5 poin)} Rancang alur kerja \textit{To-Be} yang mengeliminasi titik gesekan tersebut. Gunakan konsep \textit{Trigger-Based Auto-Capture} dan \textit{Digital Validation}. Jelaskan alurnya dalam format naratif atau diagram sederhana.

\vspace{0.2cm}
\noindent\rule{\textwidth}{0.3pt}
\vspace{0.5cm}
\noindent\rule{\textwidth}{0.3pt}
\vspace{0.5cm}
\noindent\rule{\textwidth}{0.3pt}
\vspace{0.5cm}
\noindent\rule{\textwidth}{0.3pt}
\vspace{0.5cm}
\noindent\rule{\textwidth}{0.3pt}
\vspace{0.5cm}
\noindent\rule{\textwidth}{0.3pt}

    \item \textbf{(5 poin)} Jika \textit{cycle time} berhasil direduksi dari 6 hari menjadi 1 hari, dan setiap hari keterlambatan menyebabkan potensi denda ketidakpatuhan Rp 500.000/dokumen, hitung estimasi penghematan tahunan untuk 384 laporan inspeksi/tahun (32 menara $\times$ 12 bulan).

\vspace{0.2cm}
\noindent\rule{\textwidth}{0.3pt}
\vspace{0.5cm}
\noindent\rule{\textwidth}{0.3pt}
\vspace{0.5cm}
\noindent\rule{\textwidth}{0.3pt}
\vspace{0.5cm}
\noindent\rule{\textwidth}{0.3pt}
\end{enumerate}

\vspace{0.5cm}

% === STUDI KASUS 2 ===
\begin{plnbox}[Skenario 2: Resistensi terhadap Validasi Digital]
\small
Setelah pilot implementasi alur kerja digital di Unit Distribusi PLN Regional Y, seorang Manajer Senior menolak menggunakan validasi digital dengan alasan:

\begin{quote}
\textit{``Tanda tangan basah lebih sah secara hukum dibandingkan tanda tangan elektronik. Saya tidak percaya sistem digital bisa menjamin keaslian dokumen. Lagipula, kalau server \textit{down}, semua proses berhenti total.''}
\end{quote}
\end{plnbox}

\vspace{0.3cm}

\begin{enumerate}[leftmargin=*, label=\textbf{D2.\alph*}]
    \item \textbf{(8 poin)} Tanggapi ketiga keberatan Manajer tersebut secara berurutan, dengan merujuk pada: (1) regulasi yang relevan (PP 71/2019), (2) mekanisme \textit{hash integrity} \& \textit{audit trail}, dan (3) prosedur \textit{Exception Handling/Fallback} yang telah dirumuskan dalam SOP.

\vspace{0.2cm}
\noindent\rule{\textwidth}{0.3pt}
\vspace{0.5cm}
\noindent\rule{\textwidth}{0.3pt}
\vspace{0.5cm}
\noindent\rule{\textwidth}{0.3pt}
\vspace{0.5cm}
\noindent\rule{\textwidth}{0.3pt}
\vspace{0.5cm}
\noindent\rule{\textwidth}{0.3pt}
\vspace{0.5cm}
\noindent\rule{\textwidth}{0.3pt}
\vspace{0.5cm}
\noindent\rule{\textwidth}{0.3pt}
\vspace{0.5cm}
\noindent\rule{\textwidth}{0.3pt}

    \item \textbf{(7 poin)} Dari perspektif \textit{Change Management}, klasifikasikan jenis resistensi Manajer tersebut dan usulkan \textbf{2 strategi konkret} untuk mengatasinya berdasarkan materi Bab 6.

\vspace{0.2cm}
\noindent\rule{\textwidth}{0.3pt}
\vspace{0.5cm}
\noindent\rule{\textwidth}{0.3pt}
\vspace{0.5cm}
\noindent\rule{\textwidth}{0.3pt}
\vspace{0.5cm}
\noindent\rule{\textwidth}{0.3pt}
\vspace{0.5cm}
\noindent\rule{\textwidth}{0.3pt}
\vspace{0.5cm}
\noindent\rule{\textwidth}{0.3pt}
\end{enumerate}

% ================= KUNCI JAWABAN (FASILITATOR) =================
\newpage
\begin{center}
{\LARGE\bfseries\color{red!70!black} KUNCI JAWABAN \& RUBRIK PENILAIAN}\\[0.2cm]
{\large\color{red!70!black}\textit{(Dokumen Rahasia -- Khusus Fasilitator)}}
\end{center}

\vspace{0.5cm}

\begin{answerbox}
\textbf{PENTING:} Halaman ini dan selanjutnya \textbf{WAJIB DIPISAHKAN} dari lembar soal peserta sebelum distribusi. Fasilitator dapat mencetak bagian ini secara terpisah atau menyimpannya dalam amplop tersegel.
\end{answerbox}

\vspace{0.5cm}

% --- Kunci Bagian A ---
\subsection{Kunci Jawaban Bagian A: Pilihan Ganda (30 poin)}

\renewcommand{\arraystretch}{1.4}
\begin{center}
\begin{tabularx}{0.95\textwidth}{@{} c c X c @{}}
\toprule
\textbf{No} & \textbf{Jawaban} & \textbf{Dasar/Penjelasan Singkat} & \textbf{Kriteria} \\
\midrule
1 & \textbf{B} & Optimalisasi di lingkungan hibrida = eliminasi \textit{waste} + otomasi metadata, bukan 100\% \textit{paperless} & 4.3.1 \\
2 & \textbf{B} & Data tercipta sekali di titik transaksi, mengalir tanpa duplikasi input & 4.3.1 \\
3 & \textbf{B} & Melacak satu dokumen nyata dari hulu ke hilir, mencatat hambatan & 4.3.1 \\
4 & \textbf{A} & Routing otomatis berdasarkan aturan (hierarki, nilai, risiko) & 4.3.2 \\
5 & \textbf{B} & Tanpa \textit{Immutability Lock}, konflik versi menyebabkan hilangnya \textit{Golden Record} & 4.3.2 \\
6 & \textbf{B} & SOP Kearsipan wajib mencakup Klasifikasi, Retensi, RBAC & 4.3.3 \\
7 & \textbf{D} & 8 komponen (setelah penambahan Klasifikasi, Retensi \& Akses) & 4.3.3 \\
8 & \textbf{C} & Gunakan prosedur manual terstandar, sinkronkan setelah pulih & 4.3.3 \\
9 & \textbf{B} & Persentase arsip ditemukan pada pencarian pertama & 4.3.4 \\
10 & \textbf{C} & \textit{Content}, \textit{Structure}, \textit{Context} -- trinitas integrasi kearsipan & 4.3.4 \\
\bottomrule
\end{tabularx}
\end{center}
\renewcommand{\arraystretch}{1.0}

\vspace{0.5cm}

% --- Kunci Bagian B ---
\subsection{Kunci Jawaban Bagian B: Benar/Salah (10 poin)}

\renewcommand{\arraystretch}{1.4}
\begin{center}
\begin{tabularx}{0.95\textwidth}{@{} c c X @{}}
\toprule
\textbf{No} & \textbf{Jawaban} & \textbf{Koreksi / Penjelasan} \\
\midrule
1 & \textbf{S (Salah)} & Transfer file via API hanyalah ``interkoneksi''. Integrasi kearsipan menuntut \textit{Content + Structure + Context} utuh, termasuk \textit{audit trail}. \\
2 & \textbf{B (Benar)} & Pada arsip hibrida, yang diotomasi adalah \textit{metadata tracking} dan sinkronisasi \textit{Finding Aid} digital. \\
3 & \textbf{S (Salah)} & Evaluasi bersifat deskriptif-analitis pada level prosedur bisnis, bukan audit teknis server/infrastruktur. \\
4 & \textbf{S (Salah)} & Komponen ini WAJIB; tanpanya SOP hanyalah ``SOP IT'', bukan SOP Kearsipan yang menjamin keamanan informasi dan kepatuhan JRA. \\
5 & \textbf{B (Benar)} & \textit{Gemba Walk} = observasi langsung + wawancara aktor lapangan tentang kendala harian. \\
\bottomrule
\end{tabularx}
\end{center}
\renewcommand{\arraystretch}{1.0}

\vspace{0.3cm}
\textit{Catatan: Untuk soal yang dijawab Salah oleh peserta, koreksi bernilai +1 poin bonus jika substansial. Tanpa koreksi = hanya 1 poin.}

\vspace{0.5cm}

% --- Kunci Bagian C ---
\subsection{Kunci Jawaban Bagian C: Jawaban Singkat (30 poin)}

\begin{enumerate}[leftmargin=*, itemsep=0.4cm]

\item \textbf{3 Teknik Fact-Finding (6 poin):}
\begin{itemize}[nosep]
    \item \textbf{\textit{Document Shadowing}} (2 poin): Menjadikan satu dokumen nyata sebagai objek pelacakan dari hulu ke hilir, mencatat hambatan dan durasi transisi.
    \item \textbf{Ekstraksi \textit{Log} Sistem} (2 poin): Menarik \textit{time-stamp} dari log aplikasi operasional untuk membuktikan durasi aktual vs SLA.
    \item \textbf{\textit{Gemba Walk}} (2 poin): Wawancara langsung kepada staf lapangan mengenai kendala manual harian.
\end{itemize}

\item \textbf{2 Domain Checklist Diagnostik (6 poin):}\\
Peserta boleh menyebutkan 2 domain apa saja dari 7 domain yang ada (contoh valid: Kebijakan, Tata Kelola, SDM/Budaya, Teknologi, Proses, Keamanan, Preservasi), dengan masing-masing satu contoh pertanyaan yang relevan dengan domain tersebut (masing-masing 3 poin per pasang domain \& pertanyaan).

\item \textbf{Perbedaan Trigger-Based vs SLA-Driven (6 poin):}
\begin{itemize}[nosep]
    \item \textbf{\textit{Trigger-Based Auto-Capture}} (3 poin): Status tertentu di sistem operasional \textbf{secara otomatis memicu} pembuatan draft arsip + metadata. Inisiasi oleh \textit{event sistem}.
    \item \textbf{\textit{SLA-Driven Escalation}} (3 poin): Jika validasi \textbf{melampaui batas waktu}, sistem otomatis mengeskalasi ke atasan berikutnya. Dipicu oleh \textit{pelanggaran tenggat waktu}.
\end{itemize}

\item \textbf{Mengapa Klasifikasi, Retensi \& Akses wajib (6 poin):}
\begin{itemize}[nosep]
    \item (3 poin) Komponen ini memastikan sistem tahu: ke laci virtual mana arsip dikunci, siapa yang dilarang membacanya (RBAC), dan kapan arsip boleh dimusnahkan sesuai JRA.
    \item (3 poin) Tanpa komponen ini, pencarian arsip mungkin cepat, tetapi keamanan informasi bocor, retensi tidak terkontrol, dan arsip tidak memiliki dasar hukum pemusnahan yang sah.
\end{itemize}

\item \textbf{3 Komponen Time-Cost-Risk Saving (6 poin):}
\begin{itemize}[nosep]
    \item (2 poin) \textbf{Penghematan Waktu Siklus:} Reduksi \textit{cycle time} proses kearsipan (hari $\rightarrow$ jam).
    \item (2 poin) \textbf{Reduksi Biaya Kesalahan:} Penghematan dari penurunan revisi, pencetakan ulang, denda keterlambatan.
    \item (2 poin) \textbf{Mitigasi Risiko:} Pengurangan temuan audit, sanksi ketidakpatuhan, dan biaya litigasi.
\end{itemize}

\end{enumerate}

\vspace{0.5cm}

% --- Kunci Bagian D ---
\newpage
\subsection{Kunci Jawaban Bagian D: Studi Kasus Analitis (30 poin)}

\textbf{Studi Kasus 1 -- Unit Transmisi PLN Regional X}

\begin{enumerate}[leftmargin=*, label=\textbf{D1.\alph*}]

\item \textbf{3 Titik Gesekan (5 poin):}\\
Minimal 3 dari yang berikut (poin proporsional):
\begin{itemize}[nosep]
    \item Pengantaran formulir fisik 1--2 hari (non-value-added \textit{delay})
    \item Pengetikan ulang data dari kertas ke Excel (duplikasi input manual)
    \item Siklus cetak $\rightarrow$ TTD basah $\rightarrow$ scan (konversi bolak-balik yang sia-sia)
    \item Pengisian metadata manual yang inkonsisten (risiko \textit{retrieval failure})
    \item Upload manual ke folder bersama tanpa indeks terstruktur
\end{itemize}

\item \textbf{Alur To-Be (5 poin):}\\
Jawaban berkualitas harus memuat elemen berikut:
\begin{itemize}[nosep]
    \item (2 poin) \textbf{Trigger}: Input inspeksi langsung di lapangan via perangkat mobile/tablet $\rightarrow$ status ``Selesai'' memicu auto-generate draft arsip + metadata otomatis (nama menara, tanggal, teknisi, kategori temuan)
    \item (2 poin) \textbf{Validasi Digital}: Routing otomatis ke Supervisor untuk validasi digital dengan SLA 1$\times$24 jam; eskalasi otomatis jika lewat
    \item (1 poin) \textbf{Penyimpanan}: Arsip tervalidasi langsung masuk repositori digital terindeks dengan \textit{hash integrity}
\end{itemize}

\item \textbf{Perhitungan ROI (5 poin):}
\begin{itemize}[nosep]
    \item Reduksi \textit{cycle time}: 6 hari $\rightarrow$ 1 hari = penghematan 5 hari/dokumen
    \item Potensi denda per dokumen: 5 hari $\times$ Rp 500.000 = Rp 2.500.000/dokumen
    \item Volume tahunan: 384 laporan
    \item \textbf{Estimasi penghematan:} 384 $\times$ Rp 2.500.000 = \textbf{Rp 960.000.000/tahun}
    \item (Poin penuh jika logika perhitungan benar, meskipun angka sedikit berbeda karena interpretasi)
\end{itemize}

\end{enumerate}

\vspace{0.4cm}

\textbf{Studi Kasus 2 -- Resistensi terhadap Validasi Digital}

\begin{enumerate}[leftmargin=*, label=\textbf{D2.\alph*}]

\item \textbf{Tanggapan 3 Keberatan (8 poin):}
\begin{itemize}[nosep]
    \item (3 poin) \textbf{Keberatan 1 (TTD basah lebih sah):} PP 71/2019 Pasal 14--16 secara eksplisit mengakui keabsahan tanda tangan elektronik bersertifikat. Secara hukum, TTD elektronik yang tersertifikasi memiliki kekuatan hukum yang sama dengan TTD basah.
    \item (3 poin) \textbf{Keberatan 2 (Keaslian dokumen):} Mekanisme \textit{hash integrity} + \textit{audit trail immutable} justru menjamin keaslian lebih kuat daripada TTD basah yang bisa di-fotokopi. Setiap perubahan tercatat permanen dan tidak dapat dihapus.
    \item (2 poin) \textbf{Keberatan 3 (Server \textit{down}):} SOP sudah memuat Komponen 7 ``\textit{Exception Handling \& Fallback}'' -- prosedur input manual terstandar yang langsung aktif jika sistem \textit{down} $>$4 jam, dengan sinkronisasi data saat sistem pulih.
\end{itemize}

\item \textbf{Klasifikasi Resistensi \& 2 Strategi (7 poin):}
\begin{itemize}[nosep]
    \item (2 poin) \textbf{Jenis:} Resistensi Kepercayaan (\textit{Trust-Based}) -- ketidakpercayaan terhadap validitas dokumen digital dan keamanan sistem.
    \item (2--3 poin) \textbf{Strategi 1:} Pendampingan 1-on-1 (\textit{hands-on coaching}) -- demonstrasi langsung bahwa jejak audit digital lebih kuat secara hukum, menunjukkan bukti konkret \textit{hash integrity} pada dokumen nyata.
    \item (2--3 poin) \textbf{Strategi 2:} Pilot terbatas (\textit{small win demonstration}) -- jalankan di satu sub-unit selama 30 hari, tunjukkan hasil reduksi \textit{cycle time} dan temuan zero-error ke Manajer Senior sebagai bukti empiris.
\end{itemize}

\end{enumerate}

\vspace{0.8cm}

% --- Rubrik Penilaian Rekap ---
\subsection{Rekap Rubrik Penilaian Akhir}

\renewcommand{\arraystretch}{1.3}
\begin{center}
\begin{tabularx}{0.9\textwidth}{@{} X c c @{}}
\toprule
\textbf{Rentang Nilai} & \textbf{Kategori} & \textbf{Status Kelulusan} \\
\midrule
90 -- 100 & Sangat Kompeten (\textit{Excellent}) & \textbf{LULUS} \\
75 -- 89 & Kompeten (\textit{Proficient}) & \textbf{LULUS} \\
65 -- 74 & Cukup Kompeten (\textit{Adequate}) & \textbf{LULUS (Batas Minimal)} \\
50 -- 64 & Belum Kompeten (\textit{Developing}) & \textbf{TIDAK LULUS -- Remediasi} \\
$<$ 50 & Perlu Pembinaan (\textit{Needs Improvement}) & \textbf{TIDAK LULUS -- Ulangi Modul} \\
\bottomrule
\end{tabularx}
\end{center}
\renewcommand{\arraystretch}{1.0}

\vspace{0.3cm}
\begin{plnbox}[Catatan untuk Fasilitator]
\small
\begin{itemize}[nosep]
    \item Peserta yang memperoleh nilai 50--64 wajib mengikuti sesi remediasi (pembahasan soal + pendalaman materi) sebelum mengikuti re-test.
    \item Peserta dengan nilai $<$50 disarankan untuk mengulang modul dari awal dengan pendampingan ekstra dari fasilitator.
    \item Skor Bagian D bersifat \textit{holistic scoring} -- fasilitator menilai kedalaman analisis, bukan hanya ketepatan kata kunci. Berikan kredit parsial jika logika peserta benar meskipun terminologi kurang tepat.
    \item Simpan lembar jawaban peserta sebagai bagian dari portofolio bukti pelatihan untuk keperluan audit kompetensi.
\end{itemize}
\end{plnbox}

\end{document}
